\cleardoublepage

\section{绪论}

\subsection{研究背景}

大语言模型在自然语言理解、文本生成、代码辅助和智能问答等任务中取得了较好效果，但模型规模扩大也显著提高了推理阶段的计算和显存需求。对于标准稠密 Transformer 模型，所有层参数在每次前向计算中都会参与计算，模型参数、KV cache 和中间激活共同占用 GPU 显存。当模型部署在单卡、消费级 GPU 或显存预算有限的服务器上时，显存容量往往成为制约推理吞吐和可部署性的关键因素。

混合专家模型（Mixture-of-Experts, MoE）通过稀疏激活机制在参数规模与单次前向计算量之间形成折中。MoE 层通常包含多个专家网络，每个 token 只由路由器选择其中少数专家参与计算，从而在扩大总参数量的同时控制单次计算量~\cite{shazeer2017outrageously,fedus2022switch}。以 Mixtral 等模型为代表的 MoE 大模型进一步说明，稀疏专家结构能够在增加总参数规模的同时，使单个 token 激活的参数量保持在相对可控的范围内~\cite{jiang2024mixtral}。然而，MoE 的稀疏激活并不意味着全部专家参数都可以低成本驻留在 GPU 上。对于资源受限环境，专家权重规模仍可能超过可用显存，系统需要在 GPU 显存、CPU 内存和 CPU-GPU 传输带宽之间进行权衡。

已有大模型推理系统主要围绕连续批处理、KV cache 管理、张量卸载和执行调度进行优化。例如，PagedAttention 通过分页式 KV cache 管理提升显存利用率~\cite{kwon2023pagedattention}，FlexGen 则面向单 GPU 场景设计了跨 GPU、CPU 和磁盘的卸载策略~\cite{sheng2023flexgen}。这些工作说明，在资源受限环境下，推理性能不仅取决于算子计算效率，还受到显存容量和数据搬移路径的影响。对于 MoE 模型，专家权重的按需加载进一步放大了这一问题：如果专家不能常驻 GPU，则每个 token 的路由选择会转化为动态专家访问序列，缓存未命中会引起额外的专家权重传输，并可能增加推理等待时间。

\subsection{研究问题与挑战}

本文关注的问题是：在显存受限、专家不能全部常驻 GPU 的 MoE 大模型推理环境中，如何利用真实路由访问序列，降低专家缓存未命中和 CPU-GPU 专家权重传输带来的额外开销。

该问题主要受到以下因素限制。第一，是否需要优化专家传输并不是固定结论。对于不同模型、显存预算、KV cache 预留和 CPU-GPU 带宽配置，推理瓶颈可能来自算子计算、注意力访存、KV cache 或专家权重传输。因此，方法设计前需要先判断目标部署配置是否处于专家传输受限状态（expert-transfer-bound），即专家权重传输是否成为推理等待时间的重要来源，避免在非瓶颈路径上进行无效优化。第二，MoE 专家访问由 token 级路由器动态决定，访问序列同时受到输入请求、层号、专家热度和短窗口局部性的影响。单纯依赖静态专家热度容易忽略请求内部和跨层访问模式，而只依赖最近访问记录又可能无法捕获 MoE 路由的结构性规律。第三，完整服务系统中的专家驻留、预取、传输重叠和执行调度高度耦合，直接在运行时系统中反复修改策略成本较高，也不利于隔离分析不同信号的贡献。

因此，本文采用由资源瓶颈筛查、路由事件采集、基于访问轨迹的重放分析（trace-driven replay）和在线专家缓存策略组成的研究路线。该路线参考已有分层 Roofline 建模思想筛选适合优化的专家传输受限场景，再从模型运行或离线采集的路由轨迹中抽取词元—层级粒度的专家访问序列，随后在可控的重放模拟器中比较不同缓存、卸载和预取策略。与直接修改模型结构或路由器不同，该路线不改变模型参数和输出分布，而是在推理系统层面对专家权重驻留进行优化。

\subsection{研究意义}

本文研究具有以下意义。首先，从系统部署角度看，MoE 模型的总参数量通常远大于同等激活计算量的稠密模型。如果必须将全部专家常驻 GPU，MoE 的部署门槛会显著升高；如果采用朴素按需加载，则专家传输又可能抵消稀疏计算带来的收益。因此，专家缓存与卸载优化是 MoE 大模型在资源受限硬件上实际部署中的关键问题。

其次，从方法分析角度看，MoE 推理中的专家访问并非只能视为无结构随机序列。不同工作负载下的路由序列可能包含专家热度、短窗口复用和跨层转移等局部性特征。对这些特征进行量化，有助于解释为什么某些模型或配置适合专家卸载，也有助于判断在线启发式策略与离线上界策略之间仍有多大差距。

最后，从工程实践角度看，基于访问轨迹的重放分析可以在不频繁改动完整推理框架的条件下，比较不同策略对缓存命中率、专家传输量和传输时间代理指标的影响。本文在该基础上进一步结合运行时路由事件导出链路，使离线分析结果能够与后续推理服务系统集成保持一致。

\subsection{技术路线与方案选择}

围绕上述问题，本文不以模型压缩、专家剪枝或重新训练路由器作为主要研究对象。量化和剪枝可以降低模型存储开销，但会改变模型数值表示或结构，并且需要额外讨论精度影响；重新训练路由器可能改变专家选择分布，不适合作为本文验证系统层专家驻留策略的前提。本文也不将通用请求调度作为主线，因为请求级调度主要改变 batch 组织和排队顺序，而本文关注的核心对象是词元—层级粒度的专家权重访问序列。

本文选择专家缓存、卸载与预取优化作为切入点，原因在于该方向直接对应显存容量有限与专家权重规模较大的矛盾：有限 GPU 显存无法容纳全部专家，而 CPU-GPU 传输带宽又不足以支持无代价按需换入。具体而言，本文首先参考 MoE-Lightning 中 HRM 的分层 Roofline 建模思想判断目标配置下的专家传输瓶颈~\cite{cao2025moelightning}；随后采集真实路由事件轨迹（router event trace），将每个 token 在每一层访问的专家转化为可重放的专家引用序列（expert reference sequence）；最后在统一预算下比较 LRU、静态热度、窗口预取、Belady 离线上界和 WREC 策略。通过这一流程，本文从三个方面分析专家缓存优化问题：目标配置是否存在专家传输瓶颈，路由序列中是否存在可利用的访问局部性，以及在线缓存策略与离线上界之间仍存在多大差距。

\subsection{本文主要工作}

本文围绕资源受限 MoE 大模型推理中的专家缓存与卸载优化展开，主要工作如下：

\begin{enumerate}
    \item 参考已有分层 Roofline 建模思想，构建面向本文实验配置的资源受限场景筛查流程。本文结合显存容量、KV cache 预留、专家权重大小和 CPU-GPU 带宽，对目标配置是否处于专家传输受限状态进行筛查，为后续策略设计提供硬件约束依据。
    \item 设计并实现路由事件采集与专家访问序列建模流程。本文将 MoE 推理过程中的 token、layer、expert 和 top-$k$ 选择结果统一表示为路由事件，并进一步转化为专家引用序列，为局部性统计和缓存重放提供输入。
    \item 构建基于访问轨迹的专家缓存与卸载重放模拟器。该模拟器在固定专家缓存预算下复现专家访问过程，统计缓存未命中、专家传输量、传输时间代理指标和预取浪费等指标，用于公平比较不同在线策略和离线上界。
    \item 提出 WREC 专家缓存策略。WREC 利用工作负载统计、请求内局部性和跨层专家转移信息构造在线缓存准入评分，在不修改模型路由结果的前提下降低专家换入开销。
    \item 围绕 Mixtral 路由轨迹开展实验评估。本文以 Dolly 数据集驱动的 Mixtral 路由轨迹为主要实验对象，比较 LRU、静态热度、窗口型策略、WREC 和 Belady 离线上界，并通过消融实验分析不同信号对策略收益的贡献。
\end{enumerate}

\subsection{论文结构}

本文后续章节安排如下。

第 2 章介绍本文涉及的相关背景和研究现状，包括 MoE 模型与稀疏路由机制、大模型推理系统中的显存管理、资源受限推理中的权重卸载方法，以及 MoE 专家缓存与预取相关工作。

第 3 章给出资源受限 MoE 推理场景的问题定义与建模方法，说明路由事件、专家引用序列、缓存容量、传输代价和传输时间代理指标等核心概念，并介绍如何参考已有 HRM 思想进行专家传输瓶颈筛查。

第 4 章介绍 WREC 方法设计与实现，包括路由轨迹采集流程、专家局部性统计、基线策略与离线上界设置、WREC 在线评分函数的构造以及重放模拟器的实现细节。

第 5 章给出实验设计与结果分析，包括实验环境、模型与工作负载、基于 HRM 思想的瓶颈筛查、专家局部性分析、主实验对比、消融实验、敏感性实验和局限性讨论。

第 6 章总结全文工作，归纳本文在资源受限 MoE 推理专家缓存优化方面得到的主要结论，并讨论后续在完整推理服务系统集成、更大规模 MoE 模型和解码阶段策略扩展上的改进方向。
